\section{Global poles on compact complex manifolds}
\label{sec:compact-complex-josefson}
For this section let $X$ be a connected compact complex manifold of positive dimension $n$, and choose a smooth Hermitian form $\omega>0$. Closedness is not assumed. Write $\PSH(X,\omega)$ for the locally integrable quasi-psh functions with $\ddc u\geq-\omega$, using their canonical upper semicontinuous representatives. Every global quasi-psh function becomes $\omega$-psh after multiplication by a sufficiently small positive constant. We shall pass from a local pole to a global one using bounded extremal functions. Two elementary local estimates handle the nonclosed background.

\begin{lemma}\label{lem:hermitian-bounded-products}
Finite maxima, locally bounded regularized suprema, and max pasting under the upper boundary inequality preserve $\omega$-psh functions. For bounded $u_1,\ldots,u_p\in\PSH(X,\omega)$, $1\leq p\leq n$, the mixed products of $\omega_{u_j}=\omega+\ddc u_j$ are well-defined positive currents. They have finite mass and give locally pluripolar sets zero trace mass. They are continuous under uniformly bounded decreasing sequences, and under uniformly bounded increasing sequences with the canonical upper regularized limit. The normalized family $\sup_X u=0$ is compact in $L^1(X)$ and has a uniform $L^1$ bound; its supremum is continuous under $L^1$ convergence.
\end{lemma}
\begin{proof}
For the first assertion freeze the smooth coefficients on arbitrarily small coordinate balls. There a constant closed positive form $\gamma$ can be chosen with $\omega\leq\gamma\leq\omega+\varepsilon\beta$. Add its quadratic potential and apply the ordinary psh maximum, envelope or pasting theorem. The resulting Hessian is at least $-\gamma\geq-\omega-\varepsilon\beta$. A finite small-ball cover and then $\varepsilon\downarrow0$ prove the exact lower bound $-\omega$ distributionally. The ordinary local constructions also supply the canonical upper semicontinuous representatives.

On a coordinate ball choose a smooth quadratic $q$ with $\ddc q\geq2\omega$. The $u_j+q$ are bounded psh. Expand every factor as $\ddc(u_j+q)+(\omega-\ddc q)$, defining the product by the earlier bounded products and smooth factors. A second choice of $q$ gives the same expansion by bounded multilinearity. Local radial smoothing of $u_j$ has Hessian at least $-\omega-o(1)\beta$, since convolution of the smooth coefficients of $\omega$ converges uniformly to them. Its products with background $\omega+o(1)\beta$ are positive and tend to the defined expansion. Thus the latter is positive. Bounded CLN on a finite chart cover gives finite mass; the local expansion and the zero polar mass of bounded products give the asserted trace property. Decreasing and increasing continuity follow term by term from the corresponding bounded local theorems. For an increasing sequence its pointwise and canonical limits differ only on a locally pluripolar set by \cref{cor:psh-pluripolar-discrepancy}, applied after adding $q$.

The proof of \cref{prop:normalized-global-psh-compactness} uses local psh compactness, moving maximum points and propagation of noncollapse through a finite connected cover. It applies here with these fixed quadratic corrections: their Hessians dominate $\omega$, and neither local compactness nor that propagation requires the background to be closed. Distributional convergence gives $\ddc u\geq-\omega$. The same continuous-majorant Hartogs argument proves supremum continuity. This proves all the compactness assertions, including the uniform normalized $L^1$ bound.
\end{proof}

\begin{lemma}\label{lem:hermitian-one-factor-comparison}
Let $u,v$ be bounded $\omega$-psh and let $T$ be a product of $n-1$ factors chosen among $\omega$, $\omega_u$ and $\omega_v$. Put $E=\{u<v\}$. Then
\begin{equation}\label{eq:hermitian-one-factor-comparison}
 \int_E(\omega_v-\omega_u)\wedge T
       \leq\int_E(v-u)\,\ddc T.
\end{equation}
The last derivative is a finite signed measure. Its variation is bounded by a fixed constant times the sum of the products obtained by deleting one factor from $T$ and inserting $\omega^2$, or by deleting two factors and inserting $\omega^3$.
\end{lemma}
\begin{proof}
Every factor has first derivative $\mathrm d\omega$ and second derivative $\ddc\omega$. Differentiate the local expansions from the preceding lemma. The product rule expresses $\ddc T$ as a sum of $\ddc\omega$ times products with one factor deleted and of $\mathrm d\omega\wedge\mathrm d^{\mathrm c}\omega$ times products with two deleted; signs and numerical coefficients depend only on the number of factors. On compact $X$, these smooth forms are bounded in variation by $C\omega^2$ and $C\omega^3$ when paired with positive complementary products. This proves the variation bound. For $n=1$, $T=1$ and its derivative is zero.

Set $w=v-u$. Choose smooth convex functions $F_\delta$ with $F_\delta=F_\delta'=0$ on $(-\infty,0]$, $0\leq F_\delta'\leq1$, $F_\delta''\geq0$, and $F_\delta(w)\to w^+$, $F_\delta'(w)\to\mathbf1_{\{w>0\}}$ pointwise. The bounded chain rule gives
\[
 \ddc F_\delta(w)\wedge T
   =F_\delta'(w)\ddc w\wedge T
      +F_\delta''(w)\,\mathrm dw\wedge\mathrm d^{\mathrm c}w\wedge T.
\]
Here is a justification for the nonclosed $T$. On each chart its expansion is a finite sum of closed bounded psh products times smooth forms. Each smooth form is a linear combination, with smooth bounded coefficients, of constant forms; their coefficient variations are dominated by a finite sum of positive complementary constant forms. Apply \cref{lem:bounded-psh-current-chain-rule} to those closed positive currents, and polarize for the difference $w$. Strong gradient convergence proves the displayed chain rule and all first derivative product rules after multiplication by the smooth coefficients. Bounded weighted continuity handles its Hessian terms. The gradient term is nonnegative, because $T$ is positive; this also follows by pointwise matrix positivity and the strong gradient limit.

Integration by parts on compact $X$ gives
$\int_X\ddc F_\delta(w)\wedge T=\int_XF_\delta(w)\ddc T$.
Indeed expand $\ddc(F_\delta(w)T)$, whose integral is zero; integrating the first derivative cross term once cancels it and leaves this identity. All derivative measures and products were justified in the preceding paragraph. Discard the nonnegative gradient term. Bounded convergence against the finite coefficient measures gives \eqref{eq:hermitian-one-factor-comparison}; the chosen derivatives vanish also at $w=0$, so no equality-level term is introduced.
\end{proof}

\begin{proposition}\label{prop:hermitian-small-gap-comparison}
There are constants $s_0,C>0$ depending only on $(X,\omega)$ such that, if $u,v$ are bounded $\omega$-psh, $\omega_v\geq\tfrac12\omega$ and $0<v-u\leq s<s_0$ on $E=\{u<v\}$, then
\[
 \int_E\omega_v^n\leq(1+Cs)\int_E\omega_u^n.
\]
In particular every bounded $\omega$-psh $u$ has $\int_X\omega_u^n>0$.
\end{proposition}
\begin{proof}
Put $a_k=\int_E\omega_u^k\wedge\omega^{n-k}$, $0\leq k\leq n$. Apply the preceding lemma to $T=\omega_u^k\wedge\omega^{n-k-1}$. Strict positivity of $\omega_v$ and the variation estimate give, after enlarging a fixed $B\geq1$,
\begin{equation}\label{eq:hermitian-pure-trace-recurrence}
 \tfrac12a_k\leq a_{k+1}+sB\sum_{h=0}^k a_h,
                  \qquad 0\leq k<n.
\end{equation}
In fact the error involves only $a_k,a_{k-1},a_{k-2}$, but the displayed bound suffices. Choose $s_0B(n+1)4^n\leq1/4$. Induction gives $a_k\leq4a_{k+1}$: the earlier inequalities imply $\sum_{h=0}^k a_h\leq(k+1)4^k a_k$, so the error is at most $a_k/4$. Consequently $a_k\leq4^{n-k}a_n$.

For mixed traces put $a_{j,k}=\int_E\omega_u^j\wedge\omega_v^k\wedge\omega^{n-j-k}$. Replacing one $v$ factor using \eqref{eq:hermitian-one-factor-comparison} gives
\[
 a_{j,k}\leq a_{j+1,k-1}
               +sB\sum_{p+q<j+k}a_{p,q}\qquad(k\geq1),
\]
with the same fixed enlargement of $B$. The variation bound has only smaller numbers of $u,v$ factors. Induct on $j+k$, then replace the $k$ remaining $v$ factors one by one. Since $sB\leq1$ and there are finitely many indices, this yields
$a_{j,k}\leq C_n\sum_{h=0}^{j+k}a_h\leq C'_n a_n$.
Finally telescope $\omega_v^n-\omega_u^n$ and apply the one-factor comparison to its $n$ mixed $(n-1)$-factor products. Their error traces have just been bounded by $C'_n a_n$. This proves the stated inequality.

For positivity set $m=\inf_Xu$ and $v=m+s$, with $0<s<s_0$. The strict sublevel $E=\{u<m+s\}$ is nonempty and open by upper semicontinuity, hence $\int_E\omega^n>0$. Here $\omega_v=\omega$ and $0<v-u\leq s$ on $E$. The comparison inequality forces $\int_E\omega_u^n>0$.
\end{proof}
\index{comparison principle!Hermitian small gaps}

For a set $A\subset X$ define its global extremal function and, for Borel $A$, its bounded-test capacity by
\[
 T_A=\sup\{u\in\PSH(X,\omega):u\leq0\text{ on }A\},\qquad
 c_\omega(A)=\sup_{\substack{v\in\PSH(X,\omega)\\0\leq v\leq1}}
                         \int_A\omega_v^n.
\]
The capacity is finite by local bounded CLN. A set is called globally pluripolar if it is contained in the pole set of a global quasi-psh function.

\begin{lemma}\label{lem:hermitian-extremal-pole-alternative}
Either $A$ is globally pluripolar or $T_A^*$ is a bounded nonnegative $\omega$-psh function. For a compact $K$ with nonempty interior only the second alternative occurs. Countable unions of globally pluripolar sets are globally pluripolar.
\end{lemma}
\begin{proof}
If the defining family has unbounded suprema on $X$, choose members $u_j\leq0$ on $A$ with $b_j=\sup_Xu_j\geq2^j$. The normalized $\psi_j=u_j-b_j\leq0$ have uniformly bounded $L^1$ norms. The decreasing sum $\sum_j2^{-j}\psi_j$ is locally quasi-psh with Hessian at least $-\omega$, is integrable by that bound, and is $-\infty$ on $A$. Otherwise $0\leq T_A\leq C$; local envelope regularization with a quadratic correction makes $T_A^*$ bounded and $\omega$-psh. A global pole conversely makes the suprema unbounded by adding arbitrary constants to a small rescaling of its potential.

A globally quasi-psh pole set has zero volume, since its potential is locally integrable, so it cannot contain the interior of $K$. For a countable union choose normalized $\omega$-psh pole potentials and take the same decreasing weighted sum. The uniform normalized $L^1$ bound again prevents identically infinite collapse.
\end{proof}

\begin{lemma}\label{lem:nearby-bounded-psh-traces}
Let $U,V$ be bounded psh on a Euclidean ball $B(a,r)$ and $|V-U|\leq er^2$ for $e\geq0$. Write $A=\ddc U$, $D=\ddc V$ and let $\beta$ be the Euclidean form. For $0\leq k\leq n$,
\[
 \int_{B(a,r/2)}\sum_{j=0}^k A^j\wedge D^{k-j}\wedge\beta^{n-k}
 \leq C_n\sum_{l=0}^k e^{k-l}
                         \int_{B(a,r)}A^l\wedge\beta^{n-l}.
\]
At $e=0$ the term with exponent zero has coefficient one.
\end{lemma}
\begin{proof}
Choose $(n+1)^2$ concentric intermediate radii between $r/2$ and $r$, with equal gaps, and cutoffs supported in the next ball and equal to one on the previous ball. Their Hessians have variation at most $C_nr^{-2}\beta$. If a mixed product has a $D$ factor, replace it by $A+\ddc(V-U)$ and integrate the difference by parts against the cutoff. Closedness of the remaining products bounds its absolute value by $C_ne$ times the integral on the next ball of the remaining product wedged with $\beta$. The other term has one fewer $D$ factor and the same total degree. Repeat until it has only $A$ factors, and apply the same procedure to every error term, whose total degree is smaller. A branch takes at most $k+(k-1)+\cdots+1$ steps, fewer than the chosen number of radii. Each term of final degree $l$ has acquired $k-l$ factors $e$. There are only a dimension-dependent number of branches and coefficients. Summing these bounds proves the formula. Bounded recursion justifies each integration by parts; no derivatives of $U,V$ are required.
\end{proof}

\begin{proposition}\label{prop:hermitian-extremal-mass-support}
If $K\subset X$ is compact and $u=T_K^*$ is bounded, then $\omega_u^n$ is supported on $K$. Put $S=\max\{1,\sup_Xu\}$ and $M=\int_X\omega_u^n$. Then
\begin{equation}\label{eq:hermitian-extremal-capacity-bound}
                    0<M\leq S^n c_\omega(K).
\end{equation}
\end{proposition}
\begin{proof}
We prove the support assertion using only the closed-background balayage of \cref{cor:homogeneous-psh-balayage}. Work in a chart disjoint from $K$. On $B(a,r)$ freeze the coefficients of $\omega$ at $a$, obtaining a constant form $\omega(a)=\ddc h_a$. Smoothness gives
$-\delta\beta\leq\omega-\omega(a)\leq\delta\beta$ with $\delta\leq C r$ on a fixed smaller chart. Set $q=\pi|z-a|^2$, so $\ddc q=\beta$.

Choquet selection gives increasing bounded competitors $0\leq u_j\leq u$, each at most zero on $K$, with canonical limit $u$. Balayage the bounded psh $U_j=u_j+h_a+\delta q$ in $B(a,r)$, obtaining $V_j\geq U_j$ with $(\ddc V_j)^n=0$. Put $\rho=q-\pi r^2$. The function
$W_j=V_j-h_a-\delta q+2\delta\rho$ is $\omega$-psh, because
$\omega+\ddc W_j\geq\omega-\omega(a)+\delta\beta\geq0$.
Its boundary upper limits are at most $u_j$. Paste $\max\{u_j,W_j\}$ in the ball to $u_j$ outside. The preceding Hermitian pasting rule makes this a global $\omega$-psh competitor, still zero on $K$. It is therefore at most $T_K\leq u$. Hence
\[
 U_j\leq V_j\leq u+h_a+\delta q-2\delta\rho.
\]
Balayage is increasing, so the canonical limit $V$ is bounded psh, has zero Monge--Amp\`ere measure and satisfies, with $U=u+h_a+\delta q$,
\begin{equation}\label{eq:hermitian-frozen-balayage-error}
                  0\leq V-U\leq2\pi\delta r^2.
\end{equation}
These claims use increasing bounded continuity; the upper inequality survives regularization because its right side is upper semicontinuous.

Write $A=\ddc U$, $D=\ddc V$. A cutoff $\chi$ equal to one on $B(a,r/4)$ and supported in $B(a,r/2)$, the difference-of-powers identity and bounded integration by parts give
\[
 \begin{aligned}
 \int_{B(a,r/4)}\omega_u^n
 &\leq\int\chi A^n\\
 &=\int(U-V)\ddc\chi\wedge\sum_{j=0}^{n-1}A^j\wedge D^{n-1-j}\\
 &\leq C\delta\int_{B(a,r/2)}
                   \sum_{j=0}^{n-1}A^j\wedge D^{n-1-j}\wedge\beta.
 \end{aligned}
\]
The first inequality follows from $A=\omega_u+\omega(a)-\omega+\delta\beta\geq\omega_u$; the equality uses $D^n=0$. Apply the preceding trace lemma with $e=2\pi\delta$. For small $r$ take $e\leq1$. On the fixed chart choose $C_0$ so that $\Gamma=\ddc u+C_0\beta\geq A\geq0$ for every such ball. The last expression is at most
$C\delta\sum_{l=0}^{n-1}\int_{B(a,r)}\Gamma^l\wedge\beta^{n-l}$.

Cover any fixed compact subset of the chart by balls $B(a,r/4)$ whose enlarged balls $B(a,r)$ have bounded overlap, using a Euclidean grid. The measures $\Gamma^l\wedge\beta^{n-l}$ have finite mass on a fixed larger compact set by bounded CLN. Summation bounds $\omega_u^n$ of that subset by $C r$ times this fixed finite mass. Let $r\downarrow0$. A countable compact exhaustion proves zero mass off $K$.

Finally $0\leq u/S\leq1$ is $\omega$-psh and
$\omega_{u/S}=S^{-1}\omega_u+(1-S^{-1})\omega\geq S^{-1}\omega_u$.
Positivity and support give $M\leq S^n\int_K\omega_{u/S}^n\leq S^n c_\omega(K)$. Strict positivity of $M$ is the small-gap comparison proposition.
\end{proof}
\index{extremal function!Hermitian mass support}

\begin{theorem}\label{thm:compact-complex-global-josefson}
Every locally pluripolar subset of a compact complex manifold is globally pluripolar: it is contained in $\{\Phi=-\infty\}$ for a global quasi-psh function $\Phi$.
\end{theorem}
\begin{proof}
By a countable chart cover and the countable-union lemma it suffices to treat
$E=\{\varphi=-\infty\}\cap U_1$, where $\varphi$ is psh on a coordinate ball $U$ and $U_1\Subset U$. Subtract a constant so that $\varphi\leq0$ on a neighborhood of $\overline U_1$. If $E$ is empty there is nothing to prove. Suppose it is not globally pluripolar, so $0\leq T_E^*\leq C$.

Take smooth psh $\varphi_j\downarrow\varphi$ on that neighborhood, still nonpositive there, with uniformly bounded $L^1$ norms. For each positive integer $N$ put
$K_{j,N}=\{\varphi_j\leq-N\}\cap\overline U_1$.
For all sufficiently large $j$ it has nonempty interior, because $E\ne\varnothing$ and $\varphi_j\downarrow-\infty$ at any point of $E$. Thus $u_{j,N}=T_{K_{j,N}}^*$ is bounded and nonnegative. For fixed $N$ the sets increase with $j$, so their extremals decrease to a bounded $\omega$-psh $u_N$. On $\{\varphi_j<-N\}\cap U_1$ the extremal is zero: all competitors are at most zero on a neighborhood of each such point. It follows that $u_N=0$ on $\{\varphi<-N\}\cap U_1$, and hence
\[
                    0\leq u_N\leq T_E^*\leq C.
\]
As $N$ increases, $u_N$ increases. Its canonical limit $u_\infty$ is bounded $\omega$-psh. Bounded decreasing and increasing continuity give
\[
 M_{j,N}:=\int_X\omega_{u_{j,N}}^n\longrightarrow
 M_N:=\int_X\omega_{u_N}^n,\qquad
 M_N\longrightarrow M_\infty:=\int_X\omega_{u_\infty}^n>0.
\]
The last strict inequality follows from the small-gap comparison, rather than from any uniform lower bound for all bounded potentials.

On a fixed neighborhood of $\overline U_1$, a quadratic $h$ with $\ddc h\geq\omega$ makes every global test $0\leq v\leq1$ into a uniformly bounded local psh $v+h$. Positivity yields $\omega_v^n\leq(\ddc(v+h))^n$. On $K_{j,N}$, $-\varphi_j\geq N$, so weighted CLN gives
\[
 \int_{K_{j,N}}\omega_v^n
 \leq N^{-1}\int_{\overline U_1}(-\varphi_j)
                                      (\ddc(v+h))^n
 \leq C'/N.
\]
The constant is independent of $v,j,N$, by their uniform bounded and $L^1$ norms. Consequently $c_\omega(K_{j,N})\leq C'/N$. The extremal capacity bound implies
\[
 M_{j,N}\leq \max\{1,\sup_Xu_{j,N}\}^n C'/N.
\]
For fixed $N$, supremum continuity under decreasing bounded convergence lets $j\to\infty$. Since $u_N\leq C$, this gives $M_N\leq\max\{1,C\}^n C'/N$. Now $N\to\infty$ forces $M_\infty=0$, the contradiction.

This proves every local piece globally pluripolar. The countable-union construction gives the required $\Phi$. On a disconnected compact manifold apply the argument to its finitely many components and join the resulting potentials. In dimension zero locally pluripolar sets are empty and the constant zero suffices.
\end{proof}
\index{Josefson's theorem!compact complex manifolds}
